\documentclass[10pt]{article}
\usepackage[T1]{fontenc}
\usepackage[utf8]{inputenc}
\usepackage{amssymb}
\usepackage[a4paper,top=0.9in,bottom=0.9in,left=1.05in,right=1.05in]{geometry}
\usepackage{parskip}
\setlength{\parskip}{0.6\baselineskip plus 1pt}
\usepackage{enumitem}
\usepackage[hidelinks]{hyperref}
\setlist[itemize]{topsep=2pt,itemsep=2pt,leftmargin=1.4em}
\pagestyle{empty}

\begin{document}

\noindent
Lukas Molzberger \\
Independent researcher \\
\href{mailto:lukas.molzberger@gmx.de}{lukas.molzberger@gmx.de} \\
ORCID: \href{https://orcid.org/0000-0003-4326-5909}{0000-0003-4326-5909}

\vspace{0.5em}
\noindent 5 July 2026

\vspace{0.5em}
\noindent
To the Editors \\
\textit{Symmetry} (MDPI)

\vspace{0.5em}
\noindent Dear Editors,

I am pleased to submit the manuscript \textbf{``Emergent $U(1)\times SU(3)$ Gauge
Structure and Lorentzian Geometry from a Deterministic Brane-Lattice Substrate''}
for consideration as an Article in \textit{Symmetry}.

The manuscript addresses a question at the heart of the journal's scope: to what
extent can the fundamental symmetries of low-energy physics---Lorentz invariance and
internal gauge symmetry---be understood as \emph{emergent} rather than postulated?
Building on the geometric-phase (Berry / Wilczek--Zee) mechanism and the tradition of
emergent gauge fields and analogue-gravity metrics in condensed-matter systems
(Volovik; Wen; Barcel\'o--Liberati--Visser), it shows, with explicit and fully
reproducible computations, that a single deterministic elastic substrate---a
four-dimensional cubic brane lattice whose only nonlinearity is the Pythagorean link
length---can simultaneously carry:

\begin{itemize}
  \item a Lorentzian light cone whose signature $(3,1)$ is the sign pattern of the
  directional prestress, not an imposed metric; and
  \item a $U(1)\times SU(3)$ gauge structure arising as the Wilczek--Zee holonomy of a
  rank-3 carrier bundle, in which the non-abelian sector generates the \emph{full}
  Lie algebra $\mathfrak{su}(3)$ (certified Lie-closure rank~8, not merely
  $\mathfrak{so}(3)$) and the leading dynamics of both sectors are fixed to Maxwell
  and Yang--Mills form by emergent gauge invariance.
\end{itemize}

To my knowledge this is the first demonstration that both an abelian and a
non-abelian gauge sector, together with a Lorentzian cone, descend from the
\emph{same} carrier geometry of one deterministic lattice, with the
$\mathfrak{su}(3)$ content exhibited explicitly rather than asserted. In keeping with
the journal's reproducibility standards, every quantitative claim is emitted by a
named, openly available script.

The manuscript is deliberate about scope: it is presented as a constructive proof of
concept about emergent symmetry \emph{structure}, and its boundaries are stated
explicitly---the color-carrying background is a stable but higher-energy texture rather
than the ground state; the emergent Lorentz invariance is only approximate (a
quantified anisotropy that is itself the escape from the Weinberg--Witten no-go
theorem); and the matter sector is deferred to future work. I hope the referees will
regard this candour as a strength rather than a limitation.

The work is well suited to the journal's Physics section and its emergent-symmetry,
geometric-phase, and mathematical-physics readership.

As required, I confirm the following:

\begin{itemize}
  \item We confirm that neither the manuscript nor any parts of its content are
  currently under consideration for publication with or published in another journal.
  \item All authors have approved the manuscript and agree with its submission to
  \textit{Symmetry}.
\end{itemize}

The work is single-authored; I am an independent researcher. The research received no
external funding, and I declare no conflicts of interest.

Thank you for considering this submission.

\vspace{0.5em}
\noindent Sincerely,

\vspace{1em}
\noindent Lukas Molzberger

\end{document}
